% ============================================
% Supplementary Note 1: Null Models and Statistical Controls
% ============================================

\subsection*{Supplementary Note 1: Null Models and Statistical Controls}
\label{sec:suppl:note1}

A potential concern is that Dominance could arise from skewed abundance distributions rather than ecological interactions. To address this, we compared experimental Dominance values against two null models.

In Null Model 1 (abundance-weighted random selection), species survival probability is proportional to abundance in the combined parental-community pool, testing whether high-abundance species simply survive regardless of parental-community origin. In Null Model 2 (shuffled abundance), abundances are randomly permuted among species within each parental community before neutral mixing, testing whether the overall skewness structure alone drives Dominance.

We generated 500 null offspring communities per model and calculated the Dominance metric. Experimental data showed mean Dominance of 0.698 ($n = 83$). Both null models produced significantly lower values: abundance-weighted null mean = 0.476 ($p < 10^{-18}$), shuffled null mean = 0.557 ($p < 10^{-11}$; Mann-Whitney U tests). These results indicate that abundance skewness alone cannot explain the observed patterns; ecological interactions beyond simple abundance-weighted survival determine which parental community dominates (Extended Data Fig.~3).

\paragraph{\rev{Event-matched additive null model.}}
\rev{We also tested a stricter event-matched additive null model, $n_{C,\mathrm{null}} = n_A+n_B$, in which each null coalesced community was constructed from the same two parental communities as the corresponding experimental event (Supplementary Fig.~38). Under this additive null, an apparent Dominance classification can arise when the two parental-community abundance vectors have substantially different Euclidean norms: for non-overlapping parental-community vectors, $\mathrm{Sim}(A,A+B)$ and $\mathrm{Sim}(B,A+B)$ are proportional to $\|n_A\|_2$ and $\|n_B\|_2$, respectively. This norm-imbalance effect is possible in principle, especially when richness is low and vector norms are sensitive to abundance unevenness among a small number of taxa. In the experimental data, however, the event-matched additive null was usually classified as Mixture, including 66/83 Base medium events (79.5\%), and the observed Base medium outcomes were shifted further toward parental asymmetry than their matched additive nulls (one-sided paired Wilcoxon test, $p = 5.9 \times 10^{-14}$). Thus, the simple additive-mixture model does not explain the observed Dominance tendency.}
